% Supplementary E — Journal §V (Experimental setup)
% Full dataset descriptors including HIDSAG mineralogical subsets.

\documentclass[11pt,a4paper]{article}
% Shared preamble for all supplementary chapter documents.
% Used by each supplementary/conference/suppl_*.tex and
% supplementary/journal/suppl_*.tex via % Shared preamble for all supplementary chapter documents.
% Used by each supplementary/conference/suppl_*.tex and
% supplementary/journal/suppl_*.tex via % Shared preamble for all supplementary chapter documents.
% Used by each supplementary/conference/suppl_*.tex and
% supplementary/journal/suppl_*.tex via \input{../preamble.tex}.

\usepackage[utf8]{inputenc}
\usepackage[T1]{fontenc}
\usepackage[a4paper,margin=2.2cm]{geometry}
\usepackage{amsmath,amssymb,amsfonts,amsthm}
\usepackage{mathtools}
\usepackage{graphicx}
\usepackage{booktabs}
\usepackage{array}
\usepackage{multirow}
\usepackage{longtable}
\usepackage{algorithm}
\usepackage{algpseudocode}
\usepackage{xcolor}
\usepackage{url}
\usepackage{cite}
\usepackage[hidelinks]{hyperref}

\graphicspath{{../../figures/}}

\renewcommand{\thesection}{\Alph{section}.\arabic{section}}
\renewcommand{\thesubsection}{\thesection.\arabic{subsection}}

\theoremstyle{plain}
\newtheorem{theorem}{Theorem}[section]
\newtheorem{lemma}[theorem]{Lemma}
\newtheorem{proposition}[theorem]{Proposition}
\theoremstyle{definition}
\newtheorem{definition}[theorem]{Definition}
\theoremstyle{remark}
\newtheorem{remark}[theorem]{Remark}

\setcounter{secnumdepth}{3}
\setcounter{tocdepth}{2}

\newcommand{\R}{\mathbb{R}}
\newcommand{\E}{\mathbb{E}}
\newcommand{\KL}[2]{\mathrm{KL}\!\left(#1 \,\middle\|\, #2\right)}
\newcommand{\ari}{\mathrm{ARI}}
\newcommand{\nmi}{\mathrm{NMI}}
\newcommand{\softmax}{\mathrm{softmax}}
.

\usepackage[utf8]{inputenc}
\usepackage[T1]{fontenc}
\usepackage[a4paper,margin=2.2cm]{geometry}
\usepackage{amsmath,amssymb,amsfonts,amsthm}
\usepackage{mathtools}
\usepackage{graphicx}
\usepackage{booktabs}
\usepackage{array}
\usepackage{multirow}
\usepackage{longtable}
\usepackage{algorithm}
\usepackage{algpseudocode}
\usepackage{xcolor}
\usepackage{url}
\usepackage{cite}
\usepackage[hidelinks]{hyperref}

\graphicspath{{../../figures/}}

\renewcommand{\thesection}{\Alph{section}.\arabic{section}}
\renewcommand{\thesubsection}{\thesection.\arabic{subsection}}

\theoremstyle{plain}
\newtheorem{theorem}{Theorem}[section]
\newtheorem{lemma}[theorem]{Lemma}
\newtheorem{proposition}[theorem]{Proposition}
\theoremstyle{definition}
\newtheorem{definition}[theorem]{Definition}
\theoremstyle{remark}
\newtheorem{remark}[theorem]{Remark}

\setcounter{secnumdepth}{3}
\setcounter{tocdepth}{2}

\newcommand{\R}{\mathbb{R}}
\newcommand{\E}{\mathbb{E}}
\newcommand{\KL}[2]{\mathrm{KL}\!\left(#1 \,\middle\|\, #2\right)}
\newcommand{\ari}{\mathrm{ARI}}
\newcommand{\nmi}{\mathrm{NMI}}
\newcommand{\softmax}{\mathrm{softmax}}
.

\usepackage[utf8]{inputenc}
\usepackage[T1]{fontenc}
\usepackage[a4paper,margin=2.2cm]{geometry}
\usepackage{amsmath,amssymb,amsfonts,amsthm}
\usepackage{mathtools}
\usepackage{graphicx}
\usepackage{booktabs}
\usepackage{array}
\usepackage{multirow}
\usepackage{longtable}
\usepackage{algorithm}
\usepackage{algpseudocode}
\usepackage{xcolor}
\usepackage{url}
\usepackage{cite}
\usepackage[hidelinks]{hyperref}

\graphicspath{{../../figures/}}

\renewcommand{\thesection}{\Alph{section}.\arabic{section}}
\renewcommand{\thesubsection}{\thesection.\arabic{subsection}}

\theoremstyle{plain}
\newtheorem{theorem}{Theorem}[section]
\newtheorem{lemma}[theorem]{Lemma}
\newtheorem{proposition}[theorem]{Proposition}
\theoremstyle{definition}
\newtheorem{definition}[theorem]{Definition}
\theoremstyle{remark}
\newtheorem{remark}[theorem]{Remark}

\setcounter{secnumdepth}{3}
\setcounter{tocdepth}{2}

\newcommand{\R}{\mathbb{R}}
\newcommand{\E}{\mathbb{E}}
\newcommand{\KL}[2]{\mathrm{KL}\!\left(#1 \,\middle\|\, #2\right)}
\newcommand{\ari}{\mathrm{ARI}}
\newcommand{\nmi}{\mathrm{NMI}}
\newcommand{\softmax}{\mathrm{softmax}}


\title{Journal Supplementary E --- Full dataset descriptors\\
{\large Companion to: \emph{Beyond Accuracy} (Section V)}}
\author{Felipe Santib\'a\~nez-Leal}
\date{2026}

\begin{document}
\maketitle

\section{Scope}
This supplement expands Section~V with the complete provenance and
per-subset technical inventory of the eleven evaluation entities:
six labelled HSI scenes (treated in conference supplement C and
summarised below for cross-reference) plus five HIDSAG mineralogical
subsets (treated here at full depth).

\section{Labelled HSI scenes}
Six standard labelled benchmarks from the UPV/EHU Hyperspectral
Remote Sensing Scenes collection: Indian Pines (AVIRIS, 200 bands,
2812 documents in the canonical fit, 12 topics), Salinas (AVIRIS,
204, 3520, 12), Salinas-A (subset of Salinas, 204, 1320, 6), Pavia
University (ROSIS-03, 103, 1980, 9; VNIR-only sensor), Kennedy Space
Center (AVIRIS-Old, 176, 2686, 12), Botswana (EO-1 Hyperion, 145,
2775, 12). See conference supplement C for per-scene sensor and
class-catalogue details. The canonical mineralogy framing of
AVIRIS-Classic data follows the USGS \emph{Spectroscopy of Rocks
and Minerals} reference of Clark~\cite{clark1999spectroscopy},
which catalogues the wavelength regions and absorption features
that the spectral library matching of \S{}F-12 exploits.

\section{HIDSAG mineralogical subsets}

\paragraph{Source.} HIDSAG is a curated industrial database of
hyperspectral point-spectrum measurements over mineralogical
samples acquired with a SWIR sensor (the
\texttt{swir\_low} modality consumed in this paper covers
$\sim 1100$--$2500$~nm in 268 contiguous narrow bands). The database
exposes five thematic subsets distributed as zip archives.

\subsection{GEOCHEM}
\textbf{Content.} 28 samples, 106 crop ROIs, 318 hyperspectral
cubes, 18 geochemical target variables: Al, As, Ca, Cu, Fe, K, Mo,
Nb, P, Pb, Rb, S, Si, Sr, Ti, Y, Zn, Zr (in ppm or wt-\%, depending
on element). Each sample carries a per-ROI cube and per-ROI
geochemistry vector; the document construction averages cube pixels
within each ROI to produce one document per (sample, crop,
measurement).

\subsection{GEOMET}
\textbf{Content.} 146 samples (the largest by sample count), 146
crops, 438 cubes, 5 geometallurgical target variables: Cu rec
(\% copper recovery), Lime cons (kg/t lime consumption), Mo rec
(\% molybdenum recovery), PH (pulp pH), WI (Bond work index in
kWh/t). The targets are continuous regression labels rather than
classification labels; the band-mask diagnostic on GEOMET uses
between-covariate-variance band ranking against these continuous
targets.

\subsection{MINERAL1}
\textbf{Content.} 99 samples, 99 crops, 297 cubes, 33 mineral
abundance target variables: Actinolite, Albite, Andesina (An40),
Anhydrite/Gypsum, Anorthite, Apatite, Biotite, Bornite, Bytownite
(An80), Calcite, Chalcopyrite, Chalcosite/Digenite, Chlorite,
Covelite, Epidote, Fe Oxides, Fe-Clays, K-Feldspar, Kaolinite,
Labradorite (An60), Molybdenite, Muscovite/Sericite, Oligoclasa
(An20), Other Amphibole, Others, Others Carbonates, Others Ti
Minerals, Pyrite, Pyroxene, Quartz, Rutile, Siderite, Tourmaline.
The list mixes feldspars (An20/40/60/80), sulfides
(Bornite/Chalcopyrite/Chalcosite/Covelite/Pyrite/Molybdenite),
carbonates (Calcite/Siderite), and hydrothermal indicators
(Tourmaline/Muscovite/Kaolinite) at $\sim$~\% mineralogy per ROI.

\subsection{MINERAL2}
\textbf{Content.} 20 samples (smallest), 20 crops, 60 cubes, 25
mineral abundance target variables: Alunite, Ankerite, Apatite,
Bornite, Calcite, Chalcocite, Chalcopyrite, Clinochlore, Enargite,
Gypsum, Hematite, Illite, Kaolin group, Molybdenite, Muscovite,
Orthoclase, Phengite, Pirit, Plagioclase, Pyrophyllite, Quartz,
Siderite, Sphalerite, Tenantite, Tetrahedrite. MINERAL2 covers a
slightly different mineralogical regime from MINERAL1, emphasising
high-sulfidation epithermal mineralogy (Alunite, Pyrophyllite,
Enargite). The very high topic-confidence MINERAL2 reports under the
band-mask diagnostic ($> 0.94$ across all four masks) is consistent
with the sharper mineralogical signatures in this subset.

\subsection{PORPHYRY}
\textbf{Content.} 28 samples, 56 crops, 168 cubes, 11 target
variables: Bt (\% biotite), Cal (\% calcite), Esp (\% specularite),
Hem (\% hematite), Kao (\% kaolinite), Mb (\% molybdenite), Mgt
(\% magnetite), Py (\% pyrite), Qz (\% quartz), [Cpy + Py + Bo]
(\% combined sulfides), group (categorical lithology). PORPHYRY is
the most operational-decision-relevant of the five subsets:
porphyry-copper deposits are the primary global Cu source and the
target variables track the alteration and sulfide-distribution
characteristics that drive resource estimation.

\section{Document construction on HIDSAG}

For each subset the corpus consists of one spectrum per (sample,
crop, cube) tuple, computed as the per-band mean over labelled
pixels within the crop ROI. The tokenisation follows the same
min-max + quantisation pipeline used on the labelled scenes
($Q + 1 = 13$ levels). The document count entering the canonical
fit on each subset is the cube count: GEOCHEM 318, GEOMET 438,
MINERAL1 297, MINERAL2 60, PORPHYRY 168.

\section{Sensor differences from labelled scenes}

HIDSAG is acquired with a point spectrometer rather than an imaging
array, so the spatial-coherence assumptions that motivate
Felzenszwalb / SLIC on the labelled scenes do not apply. The
framework axes F-6 (cross-method spatial agreement) and F-7
(spatial topic-label coupling on the H~$\times$~W raster) are
therefore not run on HIDSAG; axes F-1, F-2, F-3, F-4, F-5, F-8,
F-9, F-10, F-11, F-12 are all run, with F-5 using
between-covariate-variance band ranking instead of the labelled
Fisher discriminant.

\section{Continuous-measurement layer}

A complementary derived layer at
\texttt{data/derived/hidsag\_topic\_measurements/<subset>.json}
joins the per-document $\theta$ matrix shipped by the band-mask
builder with the continuous mineralogical / geometallurgical assay
values extracted from each sample's \texttt{metadata.json} (per the
HIDSAG figshare data release~\cite{hidsag2023figshare}). The layer
ships 427 records across the five subsets, exposing 5--33 continuous
variables per subset:

\begin{itemize}
\item GEOMET: Cu rec, Mo rec, PH, Lime cons, WI (5 vars, 146 records)
\item MINERAL1: 33 mineral abundance \% (99 records)
\item MINERAL2: 25 mineral abundance \% (20 records)
\item GEOCHEM: 18 elemental wt\%/ppm (106 records)
\item PORPHYRY: 11 mineral abundance \% incl.\ biotite / molybdenite
/ chalcopyrite (56 records)
\end{itemize}

This layer feeds the topic-conditional ridges, the GEOMET corner
plot, and the HIDSAG mosaic of supplement I. Figure
\ref{fig:hidsag-ridges} of Suppl I reports the most striking signal:
on PORPHYRY, per-topic mean \% biotite ranges $1.1$--$39.6$,
indicating that the canonical-fit topic basis on PORPHYRY captures
the porphyry-alteration mineralogy at high fidelity. The extraction
builder is documented at
\texttt{data-pipeline/build\_hidsag\_topic\_measurements.py}.

\section{Citation status}

HIDSAG is cited via two complementary bibliographic entries:

\begin{itemize}
\item \textbf{\texttt{hidsag2023}} --- Ehrenfeld, Egaña,
Santibáñez-Leal, Garrido, Ojeda, Townley, and Navarro,
``HIDSAG: Hyperspectral Image Database for Supervised Analysis in
Geometallurgy'', \emph{Scientific Data} 10:164 (2023),
doi:10.1038/s41597-023-02061-x. The canonical paper introducing the
database, its acquisition protocol, and the five thematic subsets.
\item \textbf{\texttt{hidsag2023figshare}} --- Ehrenfeld (2023),
HIDSAG figshare Collection,
doi:10.6084/m9.figshare.c.5983921.v1. The data release accompanying
the paper.
\end{itemize}

The reference implementation is at \url{https://github.com/alges/hidsag}.
The companion code repository for this manuscript consumes a curated
subset documented at
\texttt{data/derived/core/hidsag\_curated\_subset.json}. Felipe
Santibáñez-Leal is a co-author of \texttt{hidsag2023}; the present
manuscript is methodologically distinct (multi-axis evaluation
framework for topic models applied to HIDSAG) and cites HIDSAG as a
data source rather than re-publishing its content.

\bibliographystyle{IEEEtran}
\bibliography{../../bibliography/refs}

\end{document}
